\hypertarget{mp_example_example_c}{}\section{Example programs using the R\+N\+Alib C library}\label{mp_example_example_c}
\hypertarget{mp_example_examples_c_old_API}{}\subsection{Using the \textquotesingle{}old\textquotesingle{} A\+PI}\label{mp_example_examples_c_old_API}
The following program exercises most commonly used functions of the library. The program folds two sequences using both the mfe and partition function algorithms and calculates the tree edit and profile distance of the resulting structures and base pairing probabilities.

\begin{DoxyNote}{Note}
This program uses the old A\+PI of R\+N\+Alib, which is in part already marked deprecated. Please consult the \mbox{\hyperlink{newAPI}{R\+N\+Alib A\+PI v3.\+0}} page for details of what changes are necessary to port your implementation to the new A\+PI.
\end{DoxyNote}

\begin{DoxyCodeInclude}
\textcolor{preprocessor}{#include  <stdio.h>}
\textcolor{preprocessor}{#include  <stdlib.h>}
\textcolor{preprocessor}{#include  <math.h>}
\textcolor{preprocessor}{#include  <string.h>}
\textcolor{preprocessor}{#include  "\mbox{\hyperlink{utils_8h}{utils.h}}"}
\textcolor{preprocessor}{#include  "\mbox{\hyperlink{fold__vars_8h}{fold\_vars.h}}"}
\textcolor{preprocessor}{#include  "\mbox{\hyperlink{fold_8h}{fold.h}}"}
\textcolor{preprocessor}{#include  "\mbox{\hyperlink{part__func_8h}{part\_func.h}}"}
\textcolor{preprocessor}{#include  "\mbox{\hyperlink{inverse_8h}{inverse.h}}"}
\textcolor{preprocessor}{#include  "\mbox{\hyperlink{RNAstruct_8h}{RNAstruct.h}}"}
\textcolor{preprocessor}{#include  "\mbox{\hyperlink{treedist_8h}{treedist.h}}"}
\textcolor{preprocessor}{#include  "\mbox{\hyperlink{stringdist_8h}{stringdist.h}}"}
\textcolor{preprocessor}{#include  "\mbox{\hyperlink{profiledist_8h}{profiledist.h}}"}

\textcolor{keywordtype}{void} main()
\{
   \textcolor{keywordtype}{char} *seq1=\textcolor{stringliteral}{"CGCAGGGAUACCCGCG"}, *seq2=\textcolor{stringliteral}{"GCGCCCAUAGGGACGC"},
        *struct1,* struct2,* xstruc;
   \textcolor{keywordtype}{float} e1, e2, tree\_dist, string\_dist, profile\_dist, kT;
   \mbox{\hyperlink{structTree}{Tree}} *T1, *T2;
   \mbox{\hyperlink{structswString}{swString}} *S1, *S2;
   \textcolor{keywordtype}{float} *pf1, *pf2;
   \mbox{\hyperlink{group__data__structures_ga31125aeace516926bf7f251f759b6126}{FLT\_OR\_DBL}} *bppm;
   \textcolor{comment}{/* fold at 30C instead of the default 37C */}
   \mbox{\hyperlink{group__model__details_gab4b11c8d9c758430960896bc3fe82ead}{temperature}} = 30.;      \textcolor{comment}{/* must be set *before* initializing  */}

   \textcolor{comment}{/* allocate memory for structure and fold */}
   struct1 = (\textcolor{keywordtype}{char}* ) \mbox{\hyperlink{utils_8h_ad7e1e137b3bf1f7108933d302a7f0177}{space}}(\textcolor{keyword}{sizeof}(\textcolor{keywordtype}{char})*(strlen(seq1)+1));
   e1 =  \mbox{\hyperlink{group__mfe__fold__single_gaadafcb0f140795ae62e5ca027e335a9b}{fold}}(seq1, struct1);

   struct2 = (\textcolor{keywordtype}{char}* ) \mbox{\hyperlink{utils_8h_ad7e1e137b3bf1f7108933d302a7f0177}{space}}(\textcolor{keyword}{sizeof}(\textcolor{keywordtype}{char})*(strlen(seq2)+1));
   e2 =  \mbox{\hyperlink{group__mfe__fold__single_gaadafcb0f140795ae62e5ca027e335a9b}{fold}}(seq2, struct2);

   \mbox{\hyperlink{group__mfe__fold__single_ga107fdfe5fd641868156bfd849f6866c7}{free\_arrays}}();     \textcolor{comment}{/* free arrays used in fold() */}

   \textcolor{comment}{/* produce tree and string representations for comparison */}
   xstruc = \mbox{\hyperlink{group__struct__utils_ga78d73cd54a068ef2812812771cdddc6f}{expand\_Full}}(struct1);
   T1 = \mbox{\hyperlink{treedist_8h_a08fe4d5afd385dce593b86eaf010c6e3}{make\_tree}}(xstruc);
   S1 = \mbox{\hyperlink{stringdist_8h_a3125991b3a403b3f89230474deb3f22e}{Make\_swString}}(xstruc);
   free(xstruc);

   xstruc = \mbox{\hyperlink{group__struct__utils_ga78d73cd54a068ef2812812771cdddc6f}{expand\_Full}}(struct2);
   T2 = \mbox{\hyperlink{treedist_8h_a08fe4d5afd385dce593b86eaf010c6e3}{make\_tree}}(xstruc);
   S2 = \mbox{\hyperlink{stringdist_8h_a3125991b3a403b3f89230474deb3f22e}{Make\_swString}}(xstruc);
   free(xstruc);

   \textcolor{comment}{/* calculate tree edit distance and aligned structures with gaps */}
   \mbox{\hyperlink{dist__vars_8h_aa03194c513af6b860e7b33e370b82bdb}{edit\_backtrack}} = 1;
   tree\_dist = \mbox{\hyperlink{treedist_8h_a3b21f1925f7071f46d93431a835217bb}{tree\_edit\_distance}}(T1, T2);
   \mbox{\hyperlink{treedist_8h_acbc1cb9bce582ea945e4a467c76a57aa}{free\_tree}}(T1); \mbox{\hyperlink{treedist_8h_acbc1cb9bce582ea945e4a467c76a57aa}{free\_tree}}(T2);
   \mbox{\hyperlink{group__struct__utils_ga1054c4477d53b31d79d4cb132100e87a}{unexpand\_aligned\_F}}(\mbox{\hyperlink{dist__vars_8h_ac1605fe3448ad0a0b809c4fb8f6a854a}{aligned\_line}});
   printf(\textcolor{stringliteral}{"%s\(\backslash\)n%s  %3.2f\(\backslash\)n"}, \mbox{\hyperlink{dist__vars_8h_ac1605fe3448ad0a0b809c4fb8f6a854a}{aligned\_line}}[0], \mbox{\hyperlink{dist__vars_8h_ac1605fe3448ad0a0b809c4fb8f6a854a}{aligned\_line}}[1], tree\_dist);

   \textcolor{comment}{/* same thing using string edit (alignment) distance */}
   string\_dist = \mbox{\hyperlink{stringdist_8h_a89e3c335ef17780576d7c0e713830db9}{string\_edit\_distance}}(S1, S2);
   free(S1); free(S2);
   printf(\textcolor{stringliteral}{"%s  mfe=%5.2f\(\backslash\)n%s  mfe=%5.2f  dist=%3.2f\(\backslash\)n"},
          \mbox{\hyperlink{dist__vars_8h_ac1605fe3448ad0a0b809c4fb8f6a854a}{aligned\_line}}[0], e1, \mbox{\hyperlink{dist__vars_8h_ac1605fe3448ad0a0b809c4fb8f6a854a}{aligned\_line}}[1], e2, string\_dist);

   \textcolor{comment}{/* for longer sequences one should also set a scaling factor for}
\textcolor{comment}{      partition function folding, e.g: */}
   kT = (\mbox{\hyperlink{group__model__details_gab4b11c8d9c758430960896bc3fe82ead}{temperature}}+273.15)*1.98717/1000.;  \textcolor{comment}{/* kT in kcal/mol */}
   \mbox{\hyperlink{group__model__details_gad3b22044065acc6dee0af68931b52cfd}{pf\_scale}} = exp(-e1/kT/strlen(seq1));

   \textcolor{comment}{/* calculate partition function and base pair probabilities */}
   e1 = \mbox{\hyperlink{group__pf__fold_gadc3db3d98742427e7001a7fd36ef28c2}{pf\_fold}}(seq1, struct1);
   \textcolor{comment}{/* get the base pair probability matrix for the previous run of pf\_fold() */}
   bppm = \mbox{\hyperlink{group__pf__fold_gac5ac7ee281aae1c5cc5898a841178073}{export\_bppm}}();
   pf1 = \mbox{\hyperlink{profiledist_8h_a3dff26e707a2a2e65a0f759caabde6e7}{Make\_bp\_profile\_bppm}}(bppm, strlen(seq1));

   e2 = \mbox{\hyperlink{group__pf__fold_gadc3db3d98742427e7001a7fd36ef28c2}{pf\_fold}}(seq2, struct2);
   \textcolor{comment}{/* get the base pair probability matrix for the previous run of pf\_fold() */}
   bppm = \mbox{\hyperlink{group__pf__fold_gac5ac7ee281aae1c5cc5898a841178073}{export\_bppm}}();
   pf2 = \mbox{\hyperlink{profiledist_8h_a3dff26e707a2a2e65a0f759caabde6e7}{Make\_bp\_profile\_bppm}}(bppm, strlen(seq2));

   \mbox{\hyperlink{group__pf__fold_gae73db3f49a94f0f72e067ecd12681dbd}{free\_pf\_arrays}}();  \textcolor{comment}{/* free space allocated for pf\_fold() */}

   profile\_dist = \mbox{\hyperlink{profiledist_8h_abe75e90e00a1e5dd8862944ed53dad5d}{profile\_edit\_distance}}(pf1, pf2);
   printf(\textcolor{stringliteral}{"%s  free energy=%5.2f\(\backslash\)n%s  free energy=%5.2f  dist=%3.2f\(\backslash\)n"},
          \mbox{\hyperlink{dist__vars_8h_ac1605fe3448ad0a0b809c4fb8f6a854a}{aligned\_line}}[0], e1, \mbox{\hyperlink{dist__vars_8h_ac1605fe3448ad0a0b809c4fb8f6a854a}{aligned\_line}}[1], e2, profile\_dist);

   \mbox{\hyperlink{profiledist_8h_a9b0b84a5a45761bf42d7c835dcdb3b85}{free\_profile}}(pf1); \mbox{\hyperlink{profiledist_8h_a9b0b84a5a45761bf42d7c835dcdb3b85}{free\_profile}}(pf2);
\}
\end{DoxyCodeInclude}


In a typical Unix environment you would compile this program using\+: 
\begin{DoxyCode}
cc $\{OPENMP\_CFLAGS\} -c example.c -I$\{hpath\}
\end{DoxyCode}
 and link using 
\begin{DoxyCode}
cc $\{OPENMP\_CFLAGS\} -o example -L$\{lpath\} -lRNA -lm
\end{DoxyCode}
 where {\itshape \$\{hpath\}} and {\itshape \$\{lpath\}} point to the location of the header files and library, respectively.

\begin{DoxyNote}{Note}
As default, the R\+N\+Alib is compiled with build-\/in {\itshape Open\+MP} multithreading support. Thus, when linking your own object files to the library you have to pass the compiler specific {\itshape \$\{O\+P\+E\+N\+M\+P\+\_\+\+C\+F\+L\+A\+GS\}} (e.\+g. \textquotesingle{}-\/fopenmp\textquotesingle{} for {\bfseries gcc}) even if your code does not use openmp specific code. However, in that case the {\itshape Open\+MP} flags may be ommited when compiling example.\+c
\end{DoxyNote}
\hypertarget{mp_example_examples_c_new_API}{}\subsection{Using the \textquotesingle{}new\textquotesingle{} v3.\+0 A\+PI}\label{mp_example_examples_c_new_API}
The following is a simple program that computes the M\+FE, partition function, and centroid structure by making use of the v3.\+0 A\+PI. 
\begin{DoxyCodeInclude}
\textcolor{preprocessor}{#include <stdio.h>}
\textcolor{preprocessor}{#include <stdlib.h>}
\textcolor{preprocessor}{#include <string.h>}

\textcolor{preprocessor}{#include  <\mbox{\hyperlink{data__structures_8h}{ViennaRNA/data\_structures.h}}>}
\textcolor{preprocessor}{#include  <\mbox{\hyperlink{params_8h}{ViennaRNA/params.h}}>}
\textcolor{preprocessor}{#include  <\mbox{\hyperlink{utils_8h}{ViennaRNA/utils.h}}>}
\textcolor{preprocessor}{#include  <\mbox{\hyperlink{eval_8h}{ViennaRNA/eval.h}}>}
\textcolor{preprocessor}{#include  <\mbox{\hyperlink{fold_8h}{ViennaRNA/fold.h}}>}
\textcolor{preprocessor}{#include  <\mbox{\hyperlink{part__func_8h}{ViennaRNA/part\_func.h}}>}


\textcolor{keywordtype}{int} main(\textcolor{keywordtype}{int} argc, \textcolor{keywordtype}{char} *argv[])\{

  \textcolor{keywordtype}{char}  *seq = \textcolor{stringliteral}{"
      AGACGACAAGGUUGAAUCGCACCCACAGUCUAUGAGUCGGUGACAACAUUACGAAAGGCUGUAAAAUCAAUUAUUCACCACAGGGGGCCCCCGUGUCUAG"};
  \textcolor{keywordtype}{char}  *mfe\_structure = \mbox{\hyperlink{group__utils_gaf37a0979367c977edfb9da6614eebe99}{vrna\_alloc}}(\textcolor{keyword}{sizeof}(\textcolor{keywordtype}{char}) * (strlen(seq) + 1));
  \textcolor{keywordtype}{char}  *prob\_string   = \mbox{\hyperlink{group__utils_gaf37a0979367c977edfb9da6614eebe99}{vrna\_alloc}}(\textcolor{keyword}{sizeof}(\textcolor{keywordtype}{char}) * (strlen(seq) + 1));

  \textcolor{comment}{/* get a vrna\_fold\_compound with default settings */}
  \mbox{\hyperlink{group__fold__compound_structvrna__fc__s}{vrna\_fold\_compound\_t}} *vc = \mbox{\hyperlink{group__fold__compound_ga6601d994ba32b11511b36f68b08403be}{vrna\_fold\_compound}}(seq, NULL, 
      \mbox{\hyperlink{group__fold__compound_gacea5b7ee6181c485f36e2afa0e9089e4}{VRNA\_OPTION\_DEFAULT}});

  \textcolor{comment}{/* call MFE function */}
  \textcolor{keywordtype}{double} mfe = (double)\mbox{\hyperlink{group__mfe__fold_gabd3b147371ccf25c577f88bbbaf159fd}{vrna\_mfe}}(vc, mfe\_structure);

  printf(\textcolor{stringliteral}{"%s\(\backslash\)n%s (%6.2f)\(\backslash\)n"}, seq, mfe\_structure, mfe);

  \textcolor{comment}{/* rescale parameters for Boltzmann factors */}
  \mbox{\hyperlink{group__energy__parameters_gad607bc3a5b5da16400e2ca4ea5560233}{vrna\_exp\_params\_rescale}}(vc, &mfe);

  \textcolor{comment}{/* call PF function */}
  \mbox{\hyperlink{group__data__structures_ga31125aeace516926bf7f251f759b6126}{FLT\_OR\_DBL}} en = \mbox{\hyperlink{group__pf__fold_ga29e256d688ad221b78d37f427e0e99bc}{vrna\_pf}}(vc, prob\_string);

  \textcolor{comment}{/* print probability string and free energy of ensemble */}
  printf(\textcolor{stringliteral}{"%s (%6.2f)\(\backslash\)n"}, prob\_string, en);

  \textcolor{comment}{/* compute centroid structure */}
  \textcolor{keywordtype}{double} dist;
  \textcolor{keywordtype}{char} *cent = \mbox{\hyperlink{group__centroid__fold_ga0e64bb67e51963dc71cbd4d30b80a018}{vrna\_centroid}}(vc, &dist);

  \textcolor{comment}{/* print centroid structure, its free energy and mean distance to the ensemble */}
  printf(\textcolor{stringliteral}{"%s (%6.2f d=%6.2f)\(\backslash\)n"}, cent, \mbox{\hyperlink{group__eval_ga58f199f1438d794a265f3b27fc8ea631}{vrna\_eval\_structure}}(vc, cent), dist);

  \textcolor{comment}{/* free centroid structure */}
  free(cent);

  \textcolor{comment}{/* free pseudo dot-bracket probability string */}
  free(prob\_string);

  \textcolor{comment}{/* free mfe structure */}
  free(mfe\_structure);

  \textcolor{comment}{/* free memory occupied by vrna\_fold\_compound */}
  \mbox{\hyperlink{group__fold__compound_gadded6039d63f5d6509836e20321534ad}{vrna\_fold\_compound\_free}}(vc);

  \textcolor{keywordflow}{return} EXIT\_SUCCESS;
\}
\end{DoxyCodeInclude}
\hypertarget{mp_example_scripting_perl_examples}{}\section{Perl Examples}\label{mp_example_scripting_perl_examples}
\hypertarget{mp_example_scripting_perl_examples_flat}{}\subsection{Using the Flat Interface}\label{mp_example_scripting_perl_examples_flat}
Example 1\+: \char`\"{}\+Simple M\+F\+E prediction\char`\"{} 
\begin{DoxyCodeInclude}
#!/usr/bin/perl

use warnings;
use strict;

use RNA;

my $seq1 = "CGCAGGGAUACCCGCG";

# compute minimum free energy (mfe) and corresponding structure
my ($ss, $mfe) = RNA::fold($seq1);

# print output
printf "%s [ %6.2f ]\(\backslash\)n", $ss, $mfe;
\end{DoxyCodeInclude}
\hypertarget{mp_example_scripting_perl_examples_oo}{}\subsection{Using the Object Oriented (\+O\+O) Interface}\label{mp_example_scripting_perl_examples_oo}
The \textquotesingle{}fold\+\_\+compound\textquotesingle{} class that serves as an object oriented interface for \mbox{\hyperlink{group__fold__compound_ga1b0cef17fd40466cef5968eaeeff6166}{vrna\+\_\+fold\+\_\+compound\+\_\+t}}

Example 1\+: \char`\"{}\+Simple M\+F\+E prediction\char`\"{} 
\begin{DoxyCodeInclude}
#!/usr/bin/perl

use warnings;
use strict;

use RNA;

my $seq1 = "CGCAGGGAUACCCGCG";

# create new fold\_compound object
my $fc = new RNA::fold\_compound($seq1);

# compute minimum free energy (mfe) and corresponding structure
my ($ss, $mfe) = $fc->mfe();

# print output
printf "%s [ %6.2f ]\(\backslash\)n", $ss, $mfe;
\end{DoxyCodeInclude}
\hypertarget{mp_example_scripting_python_examples}{}\section{Python Examples}\label{mp_example_scripting_python_examples}
\hypertarget{mp_example_scripting_python_examples_flat}{}\subsection{Using the Flat Interface}\label{mp_example_scripting_python_examples_flat}
Example 1\+: \char`\"{}\+Simple M\+F\+E prediction\char`\"{} 
\begin{DoxyCodeInclude}
\textcolor{keyword}{import} RNA

sequence = \textcolor{stringliteral}{"CUCGUCGCCUUAAUCCAGUGCGGGCGCUAGACAUCUAGUUAUCGCCGCAA"}

\textcolor{comment}{# compute minimum free energy (mfe) and corresponding structure}
(structure, mfe) = RNA.fold(sequence)

\textcolor{comment}{# print output}
\textcolor{keywordflow}{print} \textcolor{stringliteral}{"%s\(\backslash\)n%s [ %6.2f ]"} % (structure, mfe)
\end{DoxyCodeInclude}
\hypertarget{mp_example_scripting_python_examples_oo}{}\subsection{Using the Object Oriented (\+O\+O) Interface}\label{mp_example_scripting_python_examples_oo}
The \textquotesingle{}fold\+\_\+compound\textquotesingle{} class that serves as an object oriented interface for \mbox{\hyperlink{group__fold__compound_ga1b0cef17fd40466cef5968eaeeff6166}{vrna\+\_\+fold\+\_\+compound\+\_\+t}}

Example 1\+: \char`\"{}\+Simple M\+F\+E prediction\char`\"{} 
\begin{DoxyCodeInclude}
\textcolor{keyword}{import} RNA;

sequence = \textcolor{stringliteral}{"CGCAGGGAUACCCGCG"}

\textcolor{comment}{# create new fold\_compound object}
fc = RNA.fold\_compound(sequence)

\textcolor{comment}{# compute minimum free energy (mfe) and corresponding structure}
(ss, mfe) = fc.mfe()

\textcolor{comment}{# print output}
\textcolor{keywordflow}{print} \textcolor{stringliteral}{"%s [ %6.2f ]"} % (ss, mfe)
\end{DoxyCodeInclude}


Example 2\+: \char`\"{}\+Use callback while generating suboptimal structures\char`\"{} 
\begin{DoxyCodeInclude}
\textcolor{keyword}{import} RNA

sequence = \textcolor{stringliteral}{"GGGGAAAACCCC"}

\textcolor{comment}{# Set global switch for unique ML decomposition}
RNA.cvar.uniq\_ML = 1

subopt\_data = \{ \textcolor{stringliteral}{'counter'} : 1, \textcolor{stringliteral}{'sequence'} : sequence \}

\textcolor{comment}{# Print a subopt result as FASTA record}
\textcolor{keyword}{def }print\_subopt\_result(structure, energy, data):
    \textcolor{keywordflow}{if} \textcolor{keywordflow}{not} structure == \textcolor{keywordtype}{None}:
        \textcolor{keywordflow}{print} \textcolor{stringliteral}{">subopt %d"} % data[\textcolor{stringliteral}{'counter'}]
        \textcolor{keywordflow}{print} \textcolor{stringliteral}{"%s"} % data[\textcolor{stringliteral}{'sequence'}]
        \textcolor{keywordflow}{print} \textcolor{stringliteral}{"%s [%6.2f]"} % (structure, energy)
        \textcolor{comment}{# increase structure counter}
        data[\textcolor{stringliteral}{'counter'}] = data[\textcolor{stringliteral}{'counter'}] + 1

\textcolor{comment}{# Create a 'fold\_compound' for our sequence}
a = RNA.fold\_compound(sequence)

\textcolor{comment}{# Enumerate all structures 500 dacal/mol = 5 kcal/mol arround}
\textcolor{comment}{# the MFE and print each structure using the function above}
a.subopt\_cb(500, print\_subopt\_result, subopt\_data);
\end{DoxyCodeInclude}


Example 3\+: \char`\"{}\+Revert M\+F\+E to Maximum Matching using soft constraint callbacks\char`\"{} 
\begin{DoxyCodeInclude}
\textcolor{keyword}{import} RNA

seq1 = \textcolor{stringliteral}{"CUCGUCGCCUUAAUCCAGUGCGGGCGCUAGACAUCUAGUUAUCGCCGCAA"}

\textcolor{comment}{# Turn-off dangles globally}
RNA.cvar.dangles = 0

\textcolor{comment}{# Data structure that will be passed to our MaximumMatching() callback with two components:}
\textcolor{comment}{# 1. a 'dummy' fold\_compound to evaluate loop energies w/o constraints, 2. a fresh set of energy parameters}
mm\_data = \{ \textcolor{stringliteral}{'dummy'}: RNA.fold\_compound(seq1), \textcolor{stringliteral}{'params'}: RNA.param() \}

\textcolor{comment}{# Nearest Neighbor Parameter reversal functions}
revert\_NN = \{ 
    RNA.DECOMP\_PAIR\_HP:       \textcolor{keyword}{lambda} i, j, k, l, f, p: - f.eval\_hp\_loop(i, j) - 100,
    RNA.DECOMP\_PAIR\_IL:       \textcolor{keyword}{lambda} i, j, k, l, f, p: - f.eval\_int\_loop(i, j, k, l) - 100,
    RNA.DECOMP\_PAIR\_ML:       \textcolor{keyword}{lambda} i, j, k, l, f, p: - p.MLclosing - p.MLintern[0] - (j - i - k + l - 2) 
      * p.MLbase - 100,
    RNA.DECOMP\_ML\_ML\_STEM:    \textcolor{keyword}{lambda} i, j, k, l, f, p: - p.MLintern[0] - (l - k - 1) * p.MLbase,
    RNA.DECOMP\_ML\_STEM:       \textcolor{keyword}{lambda} i, j, k, l, f, p: - p.MLintern[0] - (j - i - k + l) * p.MLbase,
    RNA.DECOMP\_ML\_ML:         \textcolor{keyword}{lambda} i, j, k, l, f, p: - (j - i - k + l) * p.MLbase,
    RNA.DECOMP\_ML\_UP:         \textcolor{keyword}{lambda} i, j, k, l, f, p: - (j - i + 1) * p.MLbase,
    RNA.DECOMP\_EXT\_STEM:      \textcolor{keyword}{lambda} i, j, k, l, f, p: - f.E\_ext\_loop(k, l),
    RNA.DECOMP\_EXT\_STEM\_EXT:  \textcolor{keyword}{lambda} i, j, k, l, f, p: - f.E\_ext\_loop(i, k),
    RNA.DECOMP\_EXT\_EXT\_STEM:  \textcolor{keyword}{lambda} i, j, k, l, f, p: - f.E\_ext\_loop(l, j),
    RNA.DECOMP\_EXT\_EXT\_STEM1: \textcolor{keyword}{lambda} i, j, k, l, f, p: - f.E\_ext\_loop(l, j-1),
            \}

\textcolor{comment}{# Maximum Matching callback function (will be called by RNAlib in each decomposition step)}
\textcolor{keyword}{def }MaximumMatching(i, j, k, l, d, data):
    \textcolor{keywordflow}{return} revert\_NN[d](i, j, k, l, data[\textcolor{stringliteral}{'dummy'}], data[\textcolor{stringliteral}{'params'}])

\textcolor{comment}{# Create a 'fold\_compound' for our sequence}
fc = RNA.fold\_compound(seq1)

\textcolor{comment}{# Add maximum matching soft-constraints}
fc.sc\_add\_f(MaximumMatching)
fc.sc\_add\_data(mm\_data, \textcolor{keywordtype}{None})

\textcolor{comment}{# Call MFE algorithm}
(s, mm) = fc.mfe()

\textcolor{comment}{# print result}
\textcolor{keywordflow}{print} \textcolor{stringliteral}{"%s\(\backslash\)n%s (MM: %d)\(\backslash\)n"} %  (seq1, s, -mm)

\end{DoxyCodeInclude}
 